% methods.tex — Full proposal draft, Week 8
% Revised per grader feedback: new section structure (Experiment Design), Participants,
% Stimuli, Procedures, Attention Check, Manipulation Check subsections added.
% Updated N=28, PCA 75%/85%, training pairs=1,593.

\section{Experiment Design}

\subsection{Data Collection}

The project draws on two separate datasets, linked by a shared movie stimulus: the 1959 Hitchcock film \textit{North by Northwest} (NNW).

\paragraph{fNIRS data (ongoing collection, target N\,=\,40).} Each participant comes in for a single session and completes all tasks in the following fixed order: (1) the NNW movie-watching block; (2) CPT session~1 immediately followed by CPT-MEM session~1; (3) the n-back task; and (4) CPT session~2 immediately followed by CPT-MEM session~2. We currently have 28 participants with usable data (spanning participant IDs p01--p34). fNIRS is recorded with a NIRSport2 system using a 50-channel cap (42 standard channels plus 8 short-separation channels) that detects cortical surface activity via near-infrared light. Files are stored in SNIRF format and synced with task events via Lab Streaming Layer (LSL) triggers. The 42 standard channels cover bilateral prefrontal, lateral temporal-parietal, and superior posterior cortex, with no coverage over primary visual cortex. Behavioral data, including n-back accuracy, CPT hit rates, false alarm rates, reaction times, and CPT-MEM recognition confidence ratings, are recorded alongside the fNIRS via PsychoPy and serve as the outcome variables for the individual differences analyses.

\paragraph{fMRI data (open dataset, N\,=\,59).} A publicly available dataset from a collaborating lab provides the fMRI ground truth for model training: 59 participants who watched the same NNW film in an MRI scanner, yielding fully preprocessed BOLD time series across 122 brain regions (59 subjects $\times$ 122 ROIs $\times$ 886 timepoints), parcellated using a combined Yeo 17-network, Brainnetome, and subcortical atlas. These preprocessed files are fully denoised AFNI outputs (\texttt{errts.tproject.nii}). No further data collection or permissions are needed for this component.

Because the two groups of participants watched the same movie, their brain responses are temporally aligned even though they are completely different people. After trimming for preprocessing differences (the fMRI pipeline removes the first 3 seconds of video, so fMRI volume 1 corresponds to movie second 4), there are 886 seconds of shared signal to work with. Crucially, the fact that the participants are different individuals means the model has to learn generalizable, stimulus-driven patterns rather than anything person-specific.

\subsection{Participants}

Thirty-four participants were enrolled (participant IDs p01--p34). Six were excluded due to persistent bad optode contact or excessive motion artifacts, yielding an analytic sample of N\,=\,28 (mean age\,=\,20.4 years, SD\,=\,0.8, range\,=\,18.9--22.5; 16 male, 12 female). Racial and ethnic composition: White ($n\,=\,11$), Asian ($n\,=\,6$), Hispanic or Latino ($n\,=\,4$), Other ($n\,=\,4$), Black or African American ($n\,=\,3$). Participants were recruited from the University of Chicago undergraduate population, compensated with course credit or payment, and provided written informed consent under a protocol approved by the University of Chicago IRB.

\subsection{Stimuli}

\paragraph{Movie stimulus.} The shared naturalistic stimulus is an 886-second clip from \textit{North by Northwest} (Alfred Hitchcock, 1959). This film was selected for two reasons: (1) it is the same stimulus used by the open fMRI dataset, enabling temporal alignment of neural responses across independent participant groups; and (2) it is a commercially released narrative film with complex audiovisual content, known from the ISC literature to drive reliable, distributed neural responses across viewers \citep{hasson2010}. The shared temporal structure is the engine of the cross-modal mapping approach.

\paragraph{Cognitive task stimuli.} Both the n-back and CPT tasks use face images as the primary stimulus. Faces were selected because they reliably activate regions overlapping with the fNIRS cap coverage (lateral prefrontal, temporal-parietal junction), making them well-matched to the spatial resolution of the portable sensor. The face pool consists of male and female photographs; each task draws from this pool without repeating images across blocks within a session.

\subsection{Procedures}

The session proceeds in four stages, run without counterbalancing: (1) NNW movie, (2) CPT session~1 $\to$ CPT-MEM session~1, (3) n-back task, (4) CPT session~2 $\to$ CPT-MEM session~2.

\paragraph{Movie watching (NNW).} Participants watch the 886-second NNW clip while fNIRS is recorded. No behavioral response is required. This block provides the naturalistic calibration data used to train the aPCR model.

\paragraph{Continuous Performance Task (CPT).} The CPT is a sustained attention task in which participants watch a continuous stream of face photographs and press the spacebar whenever they see a face of a pre-specified target gender. Each CPT session lasts 600 seconds and contains 500 trials (1,200\,ms each, no gap between trials). On every trial, a face image appears inside a circle overlaid on a scene photograph (an indoor or outdoor background, shown at reduced opacity) alongside a central ``+'' symbol (fixation cross); both images remain on screen for the full 1,200\,ms. Approximately 10\,\% of faces belong to the target gender (the rare gender participants must respond to); the remaining 90\,\% belong to the non-target gender and require no response. Target gender is counterbalanced: participants with even IDs respond to male faces in session~1 and female faces in session~2; participants with odd IDs follow the reverse order. Because the task runs continuously for 10 minutes with no breaks or cues marking target trials, maintaining attention is effortful, and momentary lapses produce visibly slower, more variable reaction times \citep{kofler2013}.

\paragraph{CPT Recognition Memory Task (CPT-MEM).} Immediately after each CPT session, participants are told for the first time that their memory will be tested for what they just saw. Because participants were never warned during the CPT that a memory test was coming, any encoding that occurred was incidental, arising naturally from attention rather than deliberate memorization. Each CPT-MEM session contains up to 240 self-paced trials. On each trial, a single image appears on screen --- either a face or a scene background from the preceding CPT --- and stays visible until the participant responds (up to 20 seconds). Half the images are ``old'' (they actually appeared during the CPT) and half are ``new'' (they were never shown). For each image, participants press a key to indicate how confident they are that they have seen it before: 1\,=\,Definitely New, 2\,=\,Maybe New, 3\,=\,Maybe Old, 4\,=\,Definitely Old. The scene backgrounds that had served as the visual setting during the CPT are included alongside the faces, so both the central stimuli (faces) and peripheral context (backgrounds) are tested. CPT-MEM measures how well participants encoded the CPT stimuli despite never being instructed to remember them, and serves as a secondary behavioral outcome.

\paragraph{N-back task.} The n-back is a working memory task in which participants watch a rapid sequence of face images and judge, for each face, whether it matches the one that appeared a set number of steps earlier in the sequence. In the 1-back condition, participants decide whether the current face is the same as the immediately preceding face. In the more demanding 2-back condition, participants decide whether the current face matches the one from two trials ago. Participants press \textit{s} if the face is a match or \textit{d} if it is not. Each trial begins with a brief central fixation cross (500\,ms), followed by the face image (750\,ms), giving a total trial duration of 1,250\,ms; responses must be made while the face is on screen. The task consists of four blocks of 100 trials each, in the order 1-back, 1-back, 2-back, 2-back, for 400 trials total, with short self-paced breaks between blocks. A practice block of 8 trials (1-back) runs before the main task. The 2-back condition requires simultaneously holding more items in mind and updating the memory buffer on every trial, placing substantially heavier demand on attention and working memory than 1-back.

\subsubsection{Attention Check}

To verify that participants were engaged with the movie stimulus, self-report ratings of attention and engagement are collected immediately following the movie-watching block. These ratings serve as a manipulation check for the quality of the naturalistic calibration data and will be used as a covariate in sensitivity analyses if engagement variance is large.

\subsubsection{Manipulation Check}

To verify that the n-back task successfully manipulated working memory load, we confirm that accuracy is lower and reaction times are higher in 2-back than in 1-back blocks, using paired $t$-tests across participants. A significant load effect in behavior is a prerequisite for interpreting any neural load-dependent predictions as meaningful rather than artifactual.

\subsection{Behavioral Data}

The core claim of this project is that predicted whole-brain fMRI patterns derived from portable fNIRS recordings carry meaningful information about individual differences in attention. Testing this claim requires pairing each participant's predicted neural patterns with behavioral measures of how well they actually performed. All behavioral data are recorded trial-by-trial via PsychoPy and stored as CSV files. Three behavioral constructs are targeted, each linked to a specific neural predictor from the aPCR model.

\paragraph{Attentional stability (primary outcome).} Intraindividual RT variability (iSD-RT) --- the within-person standard deviation of correct-trial reaction times --- is computed separately for each n-back condition (1-back, 2-back) and for each CPT session. iSD-RT is preferred over mean RT or overall accuracy because it captures moment-to-moment lapses in attention rather than average performance level, and remains sensitive to individual differences even when most participants perform well on average \citep{kofler2013, macdonald2009}. The central analysis in Phase~5 will test whether participants with higher \textit{predicted} default mode network (DMN) activation during these tasks --- a neural signature of failing to suppress the default mode during cognitively demanding work --- also show higher iSD-RT, indicating more frequent attentional lapses.

\paragraph{Working memory capacity (secondary outcome).} For the n-back task, proportion-correct accuracy is computed separately for 1-back and 2-back conditions. Accuracy is expected to be near-ceiling in 1-back (limiting its sensitivity) but to show meaningful variance in 2-back, where the memory demand is substantially higher. This measure will be correlated with predicted activation in the dorsal attention network (DAN) and frontoparietal control network (CONT): participants whose predicted frontoparietal engagement is higher should perform better on the 2-back task, consistent with the established fMRI-behavior literature on working memory \citep{braun2015, cole2013}.

\paragraph{Sustained attention quality (secondary outcome).} For the CPT, signal detection theory indices (d-prime and criterion) are computed from hit rates and false alarm rates, with a 0.5-trial correction for boundary cases. d-prime reflects how reliably a participant distinguished target faces from non-target faces, independent of their general tendency to press or not press the spacebar. These indices will be correlated with predicted DAN and CONT activation, and CPT iSD-RT will also be correlated with predicted DMN activation, replicating the n-back analysis in a sustained attention context. Convergence across both tasks would strengthen the claim that the model captures a domain-general attentional mechanism.

\paragraph{Incidental encoding fidelity (tertiary outcome).} For the CPT-MEM task, recognition sensitivity is quantified as d-prime derived from the 4-point confidence ratings, collapsed to a binary old/new judgment. This measures how well participants encoded the CPT stimuli without being instructed to memorize them --- a function that depends on the same frontoparietal attentional systems targeted by the CPT. Faces and scene backgrounds are analyzed separately because they differ in how they were displayed and in the neural systems they engage. CPT-MEM performance will be examined as a secondary test of whether predicted neural activation during the CPT reflects genuine attentional engagement, rather than simple motor response behavior.

\subsection{Analytical Pipeline}

The analysis unfolds across five phases. Phases 1 through 4 are complete for the NNW movie data. Preprocessing of the cognitive task data is in progress, and the cross-task analyses (Phase 5) are pending.

\subsubsection{Phase 1: fNIRS Preprocessing}

Raw SNIRF files are processed in MATLAB using the NIRS Brain AnalyzIR Toolbox \citep{santosa2018}. The same pipeline runs for all tasks; only the task number parameter changes. The steps are: (1) task trimming using LSL triggers; (2) downsampling to 4\,Hz; (3) short-separation channel labeling (8 channels at 7--9\,mm source-detector distances flagged as scalp noise references); (4) intensity-to-optical-density conversion; (5) wavelet-based motion correction; (6) Modified Beer-Lambert Law conversion to HbO/HbR concentrations (PPF\,=\,0.1, DPF\,=\,6); (7) short-channel regression via OLS to remove scalp and cardiovascular noise \citep{goodwin2014}; and (8) quality control via visual inspection and leave-one-out ISC.

For NNW, data are resampled to 1\,Hz and trimmed to 886 timepoints to match the fMRI. After quality-based exclusion, the final analytic sample is N\,=\,28 subjects, 42 channels. For each cognitive task (CPT sessions~1 and~2, CPT-MEM sessions~1 and~2, n-back), the same preprocessing script runs with only the task identifier changed.

\subsubsection{Phase 2: fMRI ROI Extraction}

Preprocessed BOLD files are handled with nibabel \citep{brett2020} and nilearn \citep{abraham2014}. The 122-ROI atlas (2\,mm isotropic) is resampled into each subject's native BOLD space (2.5\,mm isotropic) using nearest-neighbor interpolation. For each of the 59 subjects, I identify the correct NNW session file, crop the BOLD data to the first 886 volumes (covering movie seconds 4--889), extract mean signal within each ROI, and z-score each ROI's time series across time. The result is an $[886 \times 122]$ matrix per subject, averaged across all 59 to produce a group-mean fMRI matrix as the primary training target.

\subsubsection{Phase 3: fNIRS-fMRI Spatial Correspondence Check}

Before building the predictive model, I verify that fNIRS channels actually correlate with fMRI activity at nearby brain locations. MNI coordinates for the 42 standard fNIRS channels are matched to their nearest fMRI ROI by Euclidean distance. Group-mean HbO time courses (with the first 3 timepoints dropped for temporal alignment) are correlated against group-mean BOLD at each matched channel-ROI pair. Significance is assessed with phase-randomization permutation tests (1,000 iterations using circular time-shifts) and FDR correction ($q < 0.05$). As a spatial specificity check, each fNIRS channel is also correlated against three remote control regions (primary auditory cortex, primary visual cortex, and global grey matter mean).

\subsubsection{Phase 4: Adapted Principal Component Regression (aPCR)}

This is the core analysis, following \citet{gao2025}. The model is trained with leave-one-fNIRS-subject-out cross-validation across N\,=\,28 iterations. In each fold, the 27 training fNIRS subjects are paired with all 59 fMRI subjects (1,593 training pairs), and their time series are concatenated along the time axis. PCA is run separately on the fNIRS side (retaining components explaining 75\% of variance, approximately 14 PCs, matching the threshold used in \citet{gao2025}) and the fMRI side (approximately 52 PCs at 85\% variance). OLS regression estimates a weight matrix $W$ mapping fNIRS PCs to fMRI PCs. To generate a prediction for the held-out subject, their HbO data are projected into fNIRS PC space, multiplied by $W$, and back-projected to 122 ROIs via the inverse fMRI PCA transform.

Model performance is evaluated with two metrics. The \textit{grp-grp} metric correlates the mean-across-LOO predicted fMRI against the group-mean true fMRI, per ROI. The \textit{indvdl-grp} metric does the same but for each individual held-out subject, then takes the median across subjects. Significance is tested with phase-randomization permutation tests (1,000 iterations, FDR corrected across 122 ROIs). I also compute a Proportion of Noise Ceiling (PNC\,=\,$r$ / Cronbach's $\alpha$ across fMRI subjects) to express model performance relative to the theoretical maximum given the data's reliability.

As a connectivity-level check, I compute inter-subject functional connectivity (ISFC) on both the true and predicted fMRI using BrainIAK \citep{kumar2022}. ISFC produces a $122 \times 122$ connectivity matrix per subject by correlating each subject's ROI time series against the group mean of all others, then averaging with Fisher $z$-transformation. If the correlation between the true and predicted ISFC matrices is significant, the model preserves not just local activation but the large-scale connectivity structure.

Once the NNW model is trained and validated, its weights are frozen. For each cognitive task, each subject's fNIRS HbO matrix is passed through the frozen model to produce a predicted whole-brain fMRI time series. No retraining or fine-tuning is performed.

\subsubsection{Phase 5: Individual Differences Analysis (Planned)}

This is the core social science contribution, where the project extends beyond the \citet{gao2025} framework. The analyses combine predicted neural data from the frozen model with behavioral performance data from the cognitive tasks, organized around three aims.

\paragraph{Aim 1: Neural plausibility of cross-task predictions.} For each cognitive task, I will compute ISFC on the predicted fMRI time series and compare it against the NNW benchmark ($r = 0.617$). I will also test whether the n-back predictions show the expected network pattern: higher predicted activation in frontoparietal attention networks (DAN, CONT) relative to DMN during task blocks, which is the well-established fMRI signature of working memory engagement \citep{owen2005, strangman2002}. For a more sensitive test, representational similarity analysis (RSA) will check whether the predicted neural geometry distinguishes 1-back from 2-back conditions at the pattern level.

\paragraph{Aim 2: Individual differences in attentional stability.} The primary behavioral measure is intraindividual reaction time variability (RTV), computed as the standard deviation of correct-trial RTs within each participant. RTV will be computed separately for the n-back (by load level: 1-back, 2-back) and CPT tasks. Per-subject predicted activation patterns will be correlated with RTV and with task accuracy. The central hypothesis, grounded in the default mode interference account \citep{sonugabarke2007}, is that participants with higher predicted DMN activation during n-back will show higher RTV, reflecting poorer DMN suppression. Secondary analyses will test whether predicted DAN and CONT engagement correlates positively with n-back accuracy and CPT d-prime.

\paragraph{Aim 3: Cross-task generalizability.} I will compare prediction quality across conditions that differ in cognitive load (1-back vs.\ 2-back) and task type (n-back vs.\ CPT vs.\ CPT-MEM). This comparison addresses whether the naturalistic mapping captures a universal cortical-subcortical relationship or one specific to certain cognitive demands.
